\documentclass[11pt]{article}

% ---------- Packages ----------
\usepackage[a4paper,margin=1in]{geometry}
\usepackage{amsmath,amssymb,amsthm,mathtools}
\usepackage{hyperref}
\usepackage{tikz}
\usepackage{tikz-cd}
\usetikzlibrary{arrows.meta,positioning,cd}

% ---------- Theorem environments ----------
\newtheorem{theorem}{Theorem}
\newtheorem{definition}{Definition}
\newtheorem{lemma}{Lemma}
\newtheorem{proposition}{Proposition}
\newtheorem{remark}{Remark}
\newtheorem{example}{Example}
\newtheorem{conjecture}{Conjecture}

% ---------- Commands ----------
\newcommand{\Meas}{\mathbf{Meas}}
\newcommand{\QSys}{\mathbf{QuerySys}}
\newcommand{\PSys}{\mathbf{ProbSys}}
\newcommand{\Kl}{\mathrm{Kl}}
\newcommand{\Pcal}{\mathcal{P}}
\newcommand{\Xcal}{\mathcal{X}}
\newcommand{\Ycal}{\mathcal{Y}}
\newcommand{\Ocal}{\mathcal{O}}
\newcommand{\Bcal}{\mathcal{B}}
\newcommand{\Qcal}{\mathcal{Q}}
\newcommand{\id}{\mathrm{id}}
\newcommand{\op}{\mathrm{op}}
\newcommand{\pr}{\mathrm{pr}}
\newcommand{\Ext}{\mathrm{Ext}}
\newcommand{\Forget}{\mathrm{Forget}}
\newcommand{\varlim}{\varprojlim}
\newcommand{\sharppush}{_\sharp}

\title{
The Extension--Marginalisation Adjunction:\\
A Categorical Statement of Paper~1
}

\author{Patrick S. Mahon}

\date{Working note --- \today}

\begin{document}

\maketitle

\begin{abstract}
The Observational Extension Theorem of Paper~1 has a clean categorical
formulation: the extension construction and the marginalisation functor
form an adjunction
\[
  \Ext \dashv \Forget : \PSys \rightleftharpoons \QSys
\]
between a category of \emph{query systems with compatible measures} and a
category of \emph{probability spaces with query structure}.  Under the
hypotheses of Paper~1 --- sequential upper-directedness, realizability, and
observational determination --- this adjunction is an \emph{equivalence of
categories}: both the unit and counit are natural isomorphisms.

The adjunction is the mathematical form of the program's philosophical
inversion.  The forgetful functor $\Forget$ extracts the marginal system from
a probability space; the extension functor $\Ext$ constructs the canonical
probability space from marginal data.  That $\Ext \dashv \Forget$ is an
equivalence says precisely that these two operations are inverses --- that
the global measure is fully determined by, and fully determines, the
coherent system of its marginals.

The residual assumption --- $\sigma$-algebras $\mathcal{O}_i$ on the
individual query outcome spaces --- is identified as the honest epistemic
boundary of the program: measurability is assumed only at the level of what
the observer can directly report, and the global measurable structure on the
realization space is derived.
\end{abstract}

\tableofcontents

%%%%%%%%%%%%%%%%%%%%%%%%%%%%%%%%%%%%%%%%%%%%%%%%%%%%%%%%%%%%%%%%%%%%%
\section{Motivation: relocating the assumption}
\label{sec:motivation}
%%%%%%%%%%%%%%%%%%%%%%%%%%%%%%%%%%%%%%%%%%%%%%%%%%%%%%%%%%%%%%%%%%%%%

The classical construction of probability assumes a $\sigma$-algebra and a
measure on a sample space $\Omega$ that the observer never directly accesses.
All observations are derived from this assumed global structure.  The
Observable Dynamics Program inverts this: it begins with what the observer
can directly report (query outcomes) and derives the global structure from
the coherence of those reports.

After the inversion, the assumption has not been eliminated but
\emph{relocated}:

\begin{center}
\renewcommand{\arraystretch}{1.4}
\begin{tabular}{lll}
\hline
\textbf{Level} & \textbf{What is assumed} & \textbf{Status} \\
\hline
$\sigma$-algebras $\Ocal_i$ on outcome spaces $O_i$
  & Observer's reportable distinctions
  & Assumed (residual) \\
$\sigma$-algebra $\Bcal_\Omega$ on realization space $\Omega$
  & Global measurable structure
  & \emph{Derived} from queries \\
Probability measure $P$ on $\Omega$
  & Global probability
  & \emph{Derived} from $\{\nu_i\}$ \\
Objectivity
  & Agreement of all coherent observers
  & \emph{Forced} by coherence \\
\hline
\end{tabular}
\end{center}

The residual assumption --- $\sigma$-algebras on each $O_i$ --- is
epistemologically grounded: it encodes the observer's capacity to distinguish
outcomes at query level $i$.  It does not presuppose any global structure.
The global measurable structure on $\Omega$ is the join
\[
  \Bcal_\Omega = \bigvee_{i \in I} \sigma(\pr_i)
  = \sigma\!\left(\pr_i^{-1}(A) : i \in I,\, A \in \Ocal_i\right),
\]
constituted by the coherence of the queries, not presupposed beneath them.

The adjunction $\Ext \dashv \Forget$ is the precise mathematical statement
that the second and third rows of this table are theorems.

\begin{remark}[The deeper question]
The residual assumption --- $\Ocal_i$ on each $O_i$ --- is itself a candidate
for derivation.  Can the $\sigma$-algebras on outcome spaces be derived from
something more primitive, such as the observer's finite discriminative
capacity (Boolean closure of distinguishable events)?  This is Stopping
Point~1 of the program and is not addressed here.  If resolved, it would
eliminate the residual assumption entirely, grounding the whole measurable
structure in operational discriminability.
\end{remark}

%%%%%%%%%%%%%%%%%%%%%%%%%%%%%%%%%%%%%%%%%%%%%%%%%%%%%%%%%%%%%%%%%%%%%
\section{The category $\QSys$ of query systems}
\label{sec:qsys}
%%%%%%%%%%%%%%%%%%%%%%%%%%%%%%%%%%%%%%%%%%%%%%%%%%%%%%%%%%%%%%%%%%%%%

\begin{definition}[Objects of $\QSys$]\label{def:qsys-objects}
An object of $\QSys$ is a pair $(F, \nu)$ where:
\begin{itemize}
  \item $F : I^\op \to \Meas$ is a \emph{query system}: a functor from a
    directed preorder $I^\op$ to $\Meas$, giving measurable spaces
    $\{(O_i, \Ocal_i)\}_{i \in I}$ and refinement maps
    $\{\pi_{ij} : O_j \to O_i\}_{i \leq j}$ with $\pi_{ii} = \id$ and
    $\pi_{ij} \circ \pi_{jk} = \pi_{ik}$.
  \item $\nu = \{\nu_i \in \Pcal O_i\}_{i \in I}$ is a \emph{compatible
    family} of probability measures: $(\pi_{ij})\sharppush \nu_j = \nu_i$
    for all $i \leq j$.
\end{itemize}
\end{definition}

\begin{definition}[Morphisms of $\QSys$]\label{def:qsys-morphisms}
A morphism $(F, \nu) \to (G, \mu)$ in $\QSys$ --- where $F : I^\op \to \Meas$
and $G : J^\op \to \Meas$ --- consists of:
\begin{itemize}
  \item A functor $\phi : I \to J$ (a map of index sets preserving the order).
  \item For each $i \in I$, a measurable map $f_i : O'_{\phi(i)} \to O_i$
    intertwining the refinements: for all $i \leq i'$ in $I$,
    \[
      f_i \circ \pi'_{\phi(i),\phi(i')} = \pi_{i,i'} \circ f_{i'},
    \]
    and pushing $\mu_{\phi(i)}$ forward to $\nu_i$:
    \[
      (f_i)\sharppush \mu_{\phi(i)} = \nu_i.
    \]
\end{itemize}
Composition of morphisms is composition of functors and pointwise
composition of the measurable maps.  The identity morphism has $\phi = \id_I$
and $f_i = \id_{O_i}$.
\end{definition}

\begin{remark}[Cofinality: when the construction is justified]
A morphism $(\phi, \{f_i\})$ always induces a measurable map
$h : \varlim O'_j \to \varlim O_i$ between realization spaces (see
Section~\ref{sec:functors}), and $h$ is always measure-preserving on
cylinders in the image of $\phi$.  The construction never breaks.

However, if $\phi$ is not \emph{cofinal} --- i.e.\ if there exist $j \in J$
with no $i \in I$ satisfying $j \leq \phi(i)$ --- then $h$ may fail to be
an isomorphism: parts of $\Omega'$ indexed by $J$-queries outside the image
of $\phi$ are invisible to $F$, and the extension $\Ext(\phi, \{f_i\})$ may
collapse or discard measure-theoretic content.

Cofinality is therefore not a gate that filters out illegal morphisms but
a \emph{condition under which the induced map is an isomorphism} ---
the condition that the source query system is comprehensive enough to
justify the completion it produces.  It is the morphism-level analogue of
realizability: just as realizability says the query projections
$\pr_i : \Omega \to O_i$ are surjective (no part of $\Omega$ escapes the
queries), cofinality says the image of $\phi$ is eventually dominating
(no part of $J$ escapes $I$).

\begin{center}
\renewcommand{\arraystretch}{1.3}
\begin{tabular}{lll}
\hline
\textbf{Condition} & \textbf{Level} & \textbf{Role} \\
\hline
Realizability & Within a query system
  & $\pr_i$ surjective --- queries cover $\Omega$ \\
Cofinality & Between query systems
  & $\phi(I)$ cofinal in $J$ --- source covers target \\
Both & & Extension is justified, not extrapolated \\
\hline
\end{tabular}
\end{center}
\end{remark}

\begin{remark}[Morphisms as query comparisons]
A morphism $(F, \nu) \to (G, \mu)$ says: the $G$-query system is at least as
fine as the $F$-query system, and the measures are compatible under this
comparison.  Each $G$-query $\phi(i)$ projects via $f_i$ to the $F$-query
$i$, with the marginal $\mu_{\phi(i)}$ pushing forward to $\nu_i$.  When
$\phi$ is cofinal, $G$ is genuinely finer: it sees everything $F$ sees and
more, and the extension construction reflects this faithfully.
\end{remark}

%%%%%%%%%%%%%%%%%%%%%%%%%%%%%%%%%%%%%%%%%%%%%%%%%%%%%%%%%%%%%%%%%%%%%
\section{The category $\PSys$ of probability spaces with query structure}
\label{sec:psys}
%%%%%%%%%%%%%%%%%%%%%%%%%%%%%%%%%%%%%%%%%%%%%%%%%%%%%%%%%%%%%%%%%%%%%

\begin{definition}[Objects of $\PSys$]\label{def:psys-objects}
An object of $\PSys$ is a triple $(\Omega, P, F)$ where:
\begin{itemize}
  \item $(\Omega, \Bcal_\Omega)$ is a measurable space.
  \item $P \in \Pcal\Omega$ is a probability measure.
  \item $F : I^\op \to \Meas$ is a query system as above, together with
    measurable projections $\{\pr_i : \Omega \to O_i\}_{i \in I}$ satisfying
    $\pi_{ij} \circ \pr_j = \pr_i$ for all $i \leq j$, and such that
    \[
      \Bcal_\Omega = \sigma\!\left(\pr_i^{-1}(A) : i \in I,\, A \in \Ocal_i\right).
    \]
\end{itemize}
The last condition says that $\Bcal_\Omega$ is \emph{exactly} the
$\sigma$-algebra generated by the query projections --- no more and no less.
$\Omega$ carries only the measurable structure that the queries can see.
\end{definition}

\begin{definition}[Morphisms of $\PSys$]\label{def:psys-morphisms}
A morphism $(\Omega, P, F) \to (\Omega', P', G)$ in $\PSys$ consists of:
\begin{itemize}
  \item A measurable map $h : \Omega \to \Omega'$ with $h\sharppush P = P'$.
  \item A morphism $(\phi, \{f_i\}) : (F, \Forget(\Omega, P, F)) \to
    (G, \Forget(\Omega', P', G))$ in $\QSys$ (defined below) such that
    for all $i \in I$:
    \[
      f_i \circ \pr'_{\phi(i)} \circ h = \pr_i.
    \]
\end{itemize}
That is, $h$ is a measure-preserving map that is compatible with the
query projections via the query morphism $(\phi, \{f_i\})$.
\end{definition}

\begin{remark}[The $\sigma$-algebra condition]
The condition $\Bcal_\Omega = \sigma(\pr_i^{-1}(A))$ is the categorical
form of \emph{observational determination}: $\Omega$ is fully described by
its query projections.  It rules out probability spaces that carry ``hidden''
measurable structure invisible to all queries.  In the program, this is
guaranteed by the fact that $\Omega = \varlim O_i$ is constructed as a
subspace of the product, so its $\sigma$-algebra is by definition generated
by the coordinate projections.
\end{remark}

%%%%%%%%%%%%%%%%%%%%%%%%%%%%%%%%%%%%%%%%%%%%%%%%%%%%%%%%%%%%%%%%%%%%%
\section{The functors}
\label{sec:functors}
%%%%%%%%%%%%%%%%%%%%%%%%%%%%%%%%%%%%%%%%%%%%%%%%%%%%%%%%%%%%%%%%%%%%%

\begin{definition}[The forgetful functor $\Forget$]\label{def:forget}
The functor $\Forget : \PSys \to \QSys$ is defined as follows.
\begin{itemize}
  \item \emph{On objects}: $\Forget(\Omega, P, F) = (F, \{(\pr_i)\sharppush P\})$.
    It retains the query system $F$ and replaces the global measure $P$ with
    its marginals.
  \item \emph{On morphisms}: a morphism $h : (\Omega, P, F) \to (\Omega', P', G)$
    maps to the underlying query morphism $(\phi, \{f_i\})$.
\end{itemize}
$\Forget$ is well-defined: the marginals $\{(\pr_i)\sharppush P\}$ form a
compatible family by functoriality of pushforward, and $(\phi, \{f_i\})$
is a morphism in $\QSys$ by the compatibility condition in the definition
of $\PSys$-morphisms.
\end{definition}

\begin{definition}[The extension functor $\Ext$]\label{def:ext}
The functor $\Ext : \QSys \to \PSys$ is defined as follows.
\begin{itemize}
  \item \emph{On objects}: $\Ext(F, \nu) = (\varlim O_i,\, P,\, F)$ where:
    \begin{itemize}
      \item $\varlim O_i$ is the projective limit measurable space with
        the subspace $\sigma$-algebra from $\prod_i O_i$,
      \item $P$ is the unique probability measure on $\varlim O_i$ with
        $(\pr_i)\sharppush P = \nu_i$ for all $i$, constructed by the
        Carathéodory extension of the finitely additive content defined
        on cylinder sets by the compatible family $\nu$.
    \end{itemize}
    The $\sigma$-algebra on $\varlim O_i$ is by construction exactly
    $\sigma(\pr_i^{-1}(A) : i \in I, A \in \Ocal_i)$, so the object
    condition of $\PSys$ is satisfied.
  \item \emph{On morphisms}: a morphism $(\phi, \{f_i\}) : (F, \nu) \to
    (G, \mu)$ in $\QSys$ induces a measurable map
    $h = \Ext(\phi, \{f_i\}) : \varlim O'_j \to \varlim O_i$
    via the \emph{universal property} of $\varlim O_i$.
    For each $i \in I$, define $g_i : \varlim O'_j \to O_i$ by
    \[
      g_i(\omega') = f_i\!\left(\pr'_{\phi(i)}(\omega')\right).
    \]
    The family $\{g_i\}$ is a compatible cone over the $I$-diagram: for
    $i \leq i'$,
    \[
      \pi_{i,i'}(g_{i'}(\omega'))
      = \pi_{i,i'}(f_{i'}(\pr'_{\phi(i')}(\omega')))
      = f_i(\pi'_{\phi(i),\phi(i')}(\pr'_{\phi(i')}(\omega')))
      = f_i(\pr'_{\phi(i)}(\omega'))
      = g_i(\omega'),
    \]
    using the intertwining condition on $\{f_i\}$ and compatibility of
    $\{\pr'_j\}$.  The universal property of $\varlim O_i$ then gives a
    unique measurable map $h : \varlim O'_j \to \varlim O_i$ with
    $\pr_i \circ h = g_i$ for all $i$.

    Measure-preservation $h\sharppush Q = P$ holds on cylinders: for
    $A \in \Ocal_i$,
    \[
      h\sharppush Q(\pr_i^{-1}(A))
      = Q(g_i^{-1}(A))
      = Q\!\left((\pr'_{\phi(i)})^{-1}(f_i^{-1}(A))\right)
      = \mu_{\phi(i)}(f_i^{-1}(A))
      = \nu_i(A)
      = P(\pr_i^{-1}(A)),
    \]
    using $(\pr'_{\phi(i)})\sharppush Q = \mu_{\phi(i)}$ and
    $(f_i)\sharppush \mu_{\phi(i)} = \nu_i$.  Since cylinders generate
    $\Bcal_\Omega$, the Determination Theorem gives $h\sharppush Q = P$.

    When $\phi$ is cofinal, $h$ is an isomorphism; in general $h$ is a
    well-defined measure-preserving map that may not be surjective.
\end{itemize}
\end{definition}

\begin{remark}[Existence and uniqueness in $\Ext$]
The definition of $\Ext$ on objects invokes the Observational Extension
Theorem (Paper~1, Theorem~6.1) for existence of $P$, and the Observational
Determination Theorem (Paper~1, Theorem~4.1) for uniqueness.  The extension
theorem requires the query system $(F, \nu)$ to admit common coarsenings,
sequential common refinements, and realizability; the determination theorem
requires only common coarsenings.

We include these as standing hypotheses on the objects of $\QSys$ and $\PSys$
respectively, so that $\Ext$ and $\Forget$ are well-defined on the full
categories.
\end{remark}

\begin{remark}[Morphisms go beyond Paper~1]
Paper~1 works entirely over a \emph{fixed} query system.  It proves existence
and uniqueness of the extension for a given $(F, \nu)$ but says nothing about
maps between different query systems.

The morphisms $(\phi, \{f_i\}) : (F, \nu) \to (G, \mu)$ in $\QSys$ ---
and the functoriality of $\Ext$ on them --- are new categorical
infrastructure built \emph{on top of} Paper~1's results.  The proof that
$\Ext$ is functorial (Section~\ref{sec:functors}) uses the projective limit
universal property together with Paper~1's Determination Theorem, but the
statement itself does not appear in Paper~1.

This means the adjunction $\Ext \dashv \Forget$ is a stronger claim than
Paper~1 alone: it asserts not just that each individual query system has a
canonical extension, but that these extensions are \emph{coherently
organised} into a functor between categories.  Verifying this coherence ---
that $\Ext$ preserves composition and identities of query morphisms --- is
the additional content that the categorical picture contributes.
\end{remark}

%%%%%%%%%%%%%%%%%%%%%%%%%%%%%%%%%%%%%%%%%%%%%%%%%%%%%%%%%%%%%%%%%%%%%
\section{The adjunction}
\label{sec:adjunction}
%%%%%%%%%%%%%%%%%%%%%%%%%%%%%%%%%%%%%%%%%%%%%%%%%%%%%%%%%%%%%%%%%%%%%

\begin{theorem}[Extension--marginalisation adjunction]\label{thm:adjunction}
The functors $\Ext : \QSys \to \PSys$ and $\Forget : \PSys \to \QSys$ form
an adjunction $\Ext \dashv \Forget$, witnessed by natural transformations:
\begin{align*}
  \eta &: \id_{\QSys} \Rightarrow \Forget \circ \Ext
    \qquad \text{(unit)}, \\
  \varepsilon &: \Ext \circ \Forget \Rightarrow \id_{\PSys}
    \qquad \text{(counit)},
\end{align*}
satisfying the triangle identities.  Under the standing hypotheses, both
$\eta$ and $\varepsilon$ are natural isomorphisms, so $\Ext \dashv \Forget$
is an \emph{equivalence of categories}.
\end{theorem}

\begin{proof}[Construction of the unit $\eta$]
For an object $(F, \nu) \in \QSys$, define
\[
  \eta_{(F,\nu)} : (F, \nu) \to \Forget(\Ext(F, \nu))
  = \left(F,\, \{(\pr_i)\sharppush P\}\right)
\]
to be the identity morphism on the query system $F$, together with the
identity maps $f_i = \id_{O_i}$.  This is a valid $\QSys$-morphism because
the extension $P$ satisfies $(\pr_i)\sharppush P = \nu_i$ by construction
(the Observational Extension Theorem), so the pushforward condition
$(f_i)\sharppush \mu_{\phi(i)} = \nu_i$ holds with $\mu_{\phi(i)} = \nu_i$
and $f_i = \id$.

Thus $\eta_{(F,\nu)}$ is the identity morphism in $\QSys$.  In particular
it is an isomorphism.  Naturality is immediate since all components are
identity morphisms.
\end{proof}

\begin{proof}[Construction of the counit $\varepsilon$]
For an object $(\Omega, P, F) \in \PSys$, define
\[
  \varepsilon_{(\Omega, P, F)} :
  \Ext(\Forget(\Omega, P, F)) \to (\Omega, P, F).
\]
We have $\Ext(\Forget(\Omega, P, F)) = \Ext(F, \{(\pr_i)\sharppush P\})
= (\varlim O_i, P', F)$ where $P'$ is the Carathéodory extension of the
marginals $\{(\pr_i)\sharppush P\}$.

The canonical map is the measurable map
\[
  h : \varlim O_i \to \Omega, \qquad
  (o_i)_{i \in I} \mapsto \text{(the unique } \omega \in \Omega
  \text{ with } \pr_i(\omega) = o_i \text{ for all } i\text{)},
\]
induced by the universal property of the projective limit in $\Meas$:
the projections $\pr_i : \Omega \to O_i$ form a compatible cone over the
diagram $F$, so they factor uniquely through $\varlim O_i$.

This map $h$ is an isomorphism of measurable spaces: it is measurable
(by the universal property), bijective (because $\Bcal_\Omega = \sigma(\pr_i^{-1}(A))$
means $\Omega$ is already a projective limit), and its inverse is measurable
for the same reason.

Measure-preservation $h\sharppush P' = P$ follows from:
both $P'$ and $h\sharppush P'$ agree with $P$ on all cylinder sets
$\pr_i^{-1}(A)$, and these generate $\Bcal_\Omega$, so by uniqueness
(Observational Determination) $h\sharppush P' = P$.

Thus $\varepsilon_{(\Omega,P,F)} = h$ is an isomorphism.  Naturality
follows from the universal property of projective limits.
\end{proof}

\begin{proof}[Triangle identities]
Since both $\eta$ and $\varepsilon$ are natural isomorphisms (in fact
the unit is the identity and the counit is induced by the projective limit
universal property), the triangle identities
\[
  \varepsilon_{\Ext(F,\nu)} \circ \Ext(\eta_{(F,\nu)}) = \id_{\Ext(F,\nu)},
  \qquad
  \Forget(\varepsilon_{(\Omega,P,F)}) \circ \eta_{\Forget(\Omega,P,F)}
  = \id_{\Forget(\Omega,P,F)}
\]
follow directly: the first because $\Ext$ applied to the identity morphism
$\eta_{(F,\nu)}$ is the identity on $\Ext(F,\nu)$, and $\varepsilon$ on
the result is the identity isomorphism; the second because $\Forget$ of an
isomorphism is an isomorphism and composing with the identity unit gives
the identity.
\end{proof}

\begin{theorem}[Equivalence of categories]\label{thm:equivalence}
Let $\QSys^{\mathrm{cof}}$ denote the subcategory of $\QSys$ consisting of
all objects but only \emph{cofinal} morphisms.  Under the standing hypotheses
(query systems sequentially upper-directed and realizable; probability spaces
observationally determined), $\QSys^{\mathrm{cof}}$ and $\PSys$ are equivalent
categories:
\[
  \QSys^{\mathrm{cof}} \simeq \PSys.
\]
The equivalence says: \emph{a query system with compatible measures and a
probability space with query structure carry exactly the same information,
provided query comparisons are comprehensive enough to justify the
completions they produce.}  Passing from one to the other via $\Ext$ or
$\Forget$ loses nothing.

On the full category $\QSys$ (allowing non-cofinal morphisms), $\Ext$ is
still a functor and $\Ext \dashv \Forget$ is still an adjunction, but the
unit and counit are no longer isomorphisms: non-cofinal morphisms produce
completions that discard information, and the adjunction witnesses this
asymmetry rather than an equivalence.
\end{theorem}

%%%%%%%%%%%%%%%%%%%%%%%%%%%%%%%%%%%%%%%%%%%%%%%%%%%%%%%%%%%%%%%%%%%%%
\section{What the equivalence means}
\label{sec:meaning}
%%%%%%%%%%%%%%%%%%%%%%%%%%%%%%%%%%%%%%%%%%%%%%%%%%%%%%%%%%%%%%%%%%%%%

\begin{remark}[The philosophical payoff]
The equivalence $\QSys \simeq \PSys$ is the mathematical form of the
program's central claim.  It says that the classical starting point
(a probability space with a measure) and the program's starting point
(a query system with compatible marginals) are interchangeable under the
program's hypotheses.

This does not mean the two starting points are identical --- $\Ext$ and
$\Forget$ are non-trivial constructions.  It means that \emph{no information
is lost or gained in either direction}: the global measure is fully
determined by the marginal system, and the marginal system is fully
determined by the global measure.

Objectivity --- the global measure $P$ --- is not assumed beneath the
observations but \emph{forced by their coherence}.  The equivalence is the
theorem that makes this precise.
\end{remark}

\begin{remark}[What the hypotheses are doing]
Paper~1 uses two distinct directedness conditions, each doing separate work.
The equivalence requires all four of the following:
\begin{enumerate}
  \item \emph{Common coarsenings} (lower-directedness): any two queries
    have a common coarsening.  This makes the cylinder family a
    $\pi$-system, which is the condition for the Determination Theorem
    (uniqueness of $P$ given its marginals) and for $\mu_0$ to be
    well-defined as a finitely additive content.  It is needed for
    $\Forget$ to be faithful.
  \item \emph{Sequential common refinements} (sequential upper-directedness):
    every countable family of queries has a common refinement.  This is
    strictly stronger than pairwise upper-directedness and is the analytic
    condition that forces $\sigma$-subadditivity of $\mu_0$, allowing
    Carathéodory's theorem to apply.  It is needed for $\Ext$ to be
    well-defined on objects.
  \item \emph{Realizability}: the projective limit $\Omega = \varlim O_i$
    has surjective projections $\pr_i$.  This is the condition that makes
    $\sigma$-additivity hold: vanishing marginals imply vanishing cylinder
    content.  Without it, $\sigma$-subadditivity fails and the Carathéodory
    extension may not produce a probability measure.
  \item \emph{Observational determination}: $\Bcal_\Omega$ is generated by
    the query projections.  This makes $\Forget$ faithful --- no hidden
    measurable structure on $\Omega$ escapes the query system.
\end{enumerate}

Conditions 1--3 are hypotheses on objects of $\QSys$.  Condition 4 is a
hypothesis on objects of $\PSys$ and is automatically satisfied for objects
in the image of $\Ext$.

The separation of conditions 1 and 2 matters: uniqueness (Determination) is
\emph{cheaper} than existence (Extension).  $\Forget$ is faithful under
condition~1 alone; $\Ext$ needs the stronger conditions 2 and~3.  The
adjunction is asymmetric in this sense: it is easier to forget a global
measure into its marginals than to recover a global measure from marginal
data.

If condition 3 (realizability) fails, $\Ext$ may not exist.  This is the
content of Conjecture~2: can realizability be replaced by a weaker or more
algebraic condition?
\end{remark}

\begin{remark}[Connection to Kolmogorov]
The classical Kolmogorov extension theorem gives the same equivalence for
projective systems of Polish spaces.  The difference is:
\begin{itemize}
  \item Kolmogorov uses \emph{topological} hypotheses (Polish spaces,
    inner regularity, Prokhorov tightness) to establish $\sigma$-additivity.
  \item The program uses \emph{realizability} --- a purely measure-theoretic
    condition on the projective limit --- to establish $\sigma$-additivity.
\end{itemize}
The two routes are complementary.  The program's route applies in settings
where the spaces need not be Polish (any standard measurable space with a
realizable query structure), and the Prokhorov/equicontinuity work on the
\texttt{equicontinuity-cfa-conjecture} branch is probing whether the
topological route can be adapted to close Conjecture~2.
\end{remark}

\begin{remark}[The Prokhorov route: a broader $\Ext$]
Paper~1 (Section~7) establishes a second, independent construction of $\Ext$
that does not require the marginals $\{\nu_i\}$ to be countably additive.
Instead it assumes only \emph{finitely additive} marginals satisfying a
\emph{surjective projective uniform tightness} (SPUT) condition --- a
compactness hypothesis on the outcome spaces that forces $\sigma$-additivity
to emerge rather than be assumed.

In categorical terms, this extends $\Ext$ to a broader category
$\QSys^{\mathrm{fa}}$ whose objects have only finitely additive compatible
families:
\[
  \Ext^{\mathrm{fa}} : \QSys^{\mathrm{fa}} \to \PSys.
\]
The SPUT condition plays the role that realizability plays in the
countably-additive route: it is the condition under which $\Ext^{\mathrm{fa}}$
is well-defined.  Under shared hypotheses, both routes produce the same
measure $P$.

This is directly relevant to Conjecture~2: the SPUT condition is a
topological/compactness condition on the outcome spaces $O_i$, and
identifying when it can be replaced by a purely algebraic condition on
the query system is one path toward closing the conjecture.  The
equicontinuity/CFA work on the \texttt{equicontinuity-cfa-conjecture}
branch is developing exactly this topological route.
\end{remark}

%%%%%%%%%%%%%%%%%%%%%%%%%%%%%%%%%%%%%%%%%%%%%%%%%%%%%%%%%%%%%%%%%%%%%
\section{The residual assumption and Stopping Point 1}
\label{sec:residual}
%%%%%%%%%%%%%%%%%%%%%%%%%%%%%%%%%%%%%%%%%%%%%%%%%%%%%%%%%%%%%%%%%%%%%

The equivalence $\QSys \simeq \PSys$ takes the $\sigma$-algebras $\Ocal_i$
on the outcome spaces $O_i$ as given.  These are the \emph{residual
assumption} of the program: the assumption that has been relocated from the
global space $\Omega$ to the local query outcome spaces.

The relocation is epistemologically meaningful:
\begin{itemize}
  \item The $\Ocal_i$ encode the observer's capacity to distinguish outcomes
    at query level $i$ --- a directly operational assumption.
  \item The $\sigma$-algebra $\Bcal_\Omega$ is not assumed but derived as
    the join $\bigvee_i \sigma(\pr_i)$.
  \item The measure $P$ is not assumed but derived by extension.
  \item Objectivity ($P$ is canonical, not a free choice) is a theorem.
\end{itemize}

\begin{conjecture}[Stopping Point 1 --- measurability from discriminability]
The $\sigma$-algebras $\Ocal_i$ on query outcome spaces can themselves be
derived from the observer's primitive discriminative capacity: the ability
to distinguish finite families of events, closed under finite Boolean
operations.

If true, the category $\QSys$ could be replaced by a category $\QSys^0$ of
query systems with only Boolean (finitely additive) event structure, and the
$\sigma$-algebras $\Ocal_i$ would emerge as the minimal countably additive
extension supporting coherent probabilistic reasoning.  The equivalence
$\QSys^0 \simeq \PSys$ would then ground the entire measurable and
probabilistic structure in operational discriminability alone.
\end{conjecture}

%%%%%%%%%%%%%%%%%%%%%%%%%%%%%%%%%%%%%%%%%%%%%%%%%%%%%%%%%%%%%%%%%%%%%
\section{Summary}
\label{sec:summary}
%%%%%%%%%%%%%%%%%%%%%%%%%%%%%%%%%%%%%%%%%%%%%%%%%%%%%%%%%%%%%%%%%%%%%

\begin{center}
\renewcommand{\arraystretch}{1.4}
\begin{tabular}{lll}
\hline
\textbf{Object} & \textbf{In $\QSys$} & \textbf{In $\PSys$} \\
\hline
Spaces & Outcome spaces $O_i$ (assumed) & $\Omega = \varlim O_i$ (constructed) \\
$\sigma$-algebras & $\Ocal_i$ on $O_i$ (assumed); $\Bcal_\Omega$ derived
  & $\Bcal_\Omega = \sigma(\pr_i^{-1}(A))$ (by construction) \\
Measures & Compatible family $\{\nu_i\}$ (given) & Global measure $P$ (derived) \\
Morphisms & Query maps $(\phi, \{f_i\})$ + measure intertwining & Measure-preserving maps + query compatibility \\
Equivalence condition & $\phi$ cofinal (covers target queries) & Observational determination ($\Bcal_\Omega$ generated by $\pr_i$) \\
\hline
\textbf{Functor} & \multicolumn{2}{l}{\textbf{Action}} \\
\hline
$\Ext$ & \multicolumn{2}{l}{$(F, \nu) \mapsto (\varlim O_i, P, F)$
  via Carathéodory extension} \\
$\Forget$ & \multicolumn{2}{l}{$(\Omega, P, F) \mapsto (F, \{(\pr_i)\sharppush P\})$
  via marginalisation} \\
\hline
\textbf{Natural transformation} & \multicolumn{2}{l}{\textbf{Content}} \\
\hline
Unit $\eta$ & \multicolumn{2}{l}{Identity: marginals of extension = original
  $\nu_i$ (Extension Theorem)} \\
Counit $\varepsilon$ & \multicolumn{2}{l}{Projective limit universal property:
  $\Omega \cong \varlim O_i$ (Determination Theorem)} \\
\hline
\end{tabular}
\end{center}

The adjunction $\Ext \dashv \Forget$ encodes both main theorems of Paper~1
in a single categorical statement.  The equivalence $\QSys \simeq \PSys$
is the categorical form of the program's philosophical inversion.

\end{document}
